\documentclass{article}
\usepackage[margin=1in]{geometry} % Adjust page margins to accommodate larger table
\usepackage{array}
\usepackage{booktabs}
\usepackage{xcolor}
\usepackage{tikz}
\usepackage{colortbl}
\usepackage{multirow}
\usepackage{enumitem}
\usepackage{graphicx}
\usepackage{wrapfig}
\usepackage{amsmath}
\geometry{left=1in, right=1in, top=1in, bottom=1in}

\definecolor{lightred}{RGB}{255,230,230}
\definecolor{darkred}{RGB}{139,0,0}

\newcommand{\sectionheading}[1]{%
  \multicolumn{2}{l}{%
    \cellcolor{darkred}\color{white}\bfseries\Large #1%
  }%
}



\begin{document}
\thispagestyle{empty} % Removes page number

% Yellow background block
\begin{tikzpicture}[remember picture, overlay]
    \fill[yellow] (current page.north east) rectangle ++(-6, -\paperheight);
\end{tikzpicture}

% Title Text in a straight line
\begin{tikzpicture}[remember picture, overlay]
    \node[anchor=north west, xshift=-1.2cm, yshift=0.25cm] {
        {\fontsize{25pt}{48pt}\selectfont
        \definecolor{darkred}{RGB}{139,0,0}\textcolor{red!85!black}{MAULANA ABUL KALAM AZAD}}
    };
\end{tikzpicture}
\begin{tikzpicture}[remember picture, overlay]
    \node[anchor=north west, xshift=-1.2cm, yshift=-1cm] {
        {\fontsize{25pt}{48pt}\selectfont
        \textcolor{black!70!black}{UNIVERSITY OF TECHNOLOGY}}
    };
\end{tikzpicture}

\begin{tikzpicture}[remember picture, overlay]
    % Absolute positioning
    \node[anchor=north west, xshift=6cm, yshift=-5.5cm] at (current page.north west) {
        \includegraphics[width=5cm]{makaut.png}
    };
\end{tikzpicture}

\begin{tikzpicture}[remember picture, overlay]
    \node[anchor=north west, xshift=0.9cm, yshift=-11.75cm] at (current page.north west) {
        {\fontsize{20pt}{24pt}\selectfont
        \textcolor{black}{\underline{Software Tools and Technology (Lab Assignment)}}}
    };
\end{tikzpicture}

\begin{center}

\vspace{9 cm}
\renewcommand{\arraystretch}{2.5}
\hspace*{-2.25in}
\begin{tabular}{ |p{3cm}|p{4cm}|p{5cm}| }
\hline
\multicolumn{3}{|c|}{\Large \textbf{\textit{Group 15}}} \\
\hline
NAME & ROLL NO.& DEPARTMENT \\
\hline
Shreyasi Mahato [LEAD]& 30001223051& Bachelor of Computer Application \\
\hline
Mouteer Datta& 30059223026& Bsc in Forensic Science\\
\hline
Rajib Roy& 30001223040& Bachelor of Computer Application\\
\hline
Priyanshu Sarma Sarkar& 300054623009& Bsc in Artificial Intelligence\\
\hline
Abhirup Mukherjee& 30085323008& Bsc in CyberSecurity\\
\hline
\end{tabular}

\end{center}

\newpage
\begin{center}
    \section*{Acknowledgment}
\end{center}

We are deeply grateful to our instructor, Dr. Ayan Ghosh, for his guidance and support throughout this project. His expertise and mentorship were invaluable in helping us successfully complete the tasks.\\

We would also like to express our sincere appreciation to the group members, to acknowledge the hard work and dedication of our team members:\\

- Shreyasi Mahato\\
- Mouteer Datta\\
- Rajib Roy\\
- Priyanshu Sarma Sarkar\\
- Abhirup Mukherjee\\

Each member contributed significantly to this project, and without their collaboration and commitment, this project would not have been possible. We thank everyone for their dedication and collaboration. Each member's contributions played a vital role in the smooth execution of this project.\\

Lastly, we thank our institution for providing us with the resources and the opportunity to learn and apply our knowledge in this project on Software Tools and Technology.

\vspace{2cm}
\noindent
Date: 23.09.2024
\\
\\
\\


Signature
\newpage
\begin{table}[h]
\centering
\large % Increase overall font size
\setlength{\arrayrulewidth}{1.5pt}
\renewcommand{\arraystretch}{1.3}
\begin{tabular}{>{\bfseries}l p{0.75\textwidth}}
\toprule[2.5pt]
\sectionheading{Index} \\
\midrule
1. & \cellcolor{yellow}Assignment 1 \\
   & My First LATEX Document \dotfill \pageref{fig:Shreyasi} \\
\midrule
2. & \cellcolor{yellow} Assignment 2 \\
   & Curriculum Vitae \dotfill {6} \\
\midrule
3. & \cellcolor{yellow}Assignment 3 \\
   & Calculator Program using C \dotfill {8} \\
\midrule
4. & \cellcolor{yellow}Assignment 4 \\
   & Matrix in LaTeX \dotfill {11} \\
   
\midrule
5. & \cellcolor{yellow}Assignment 5 \\
   & Mind Reader App \dotfill {14} \\
  
\bottomrule[2.5pt]
\end{tabular}

\end{table}

% Decorative element using TikZ
\begin{tikzpicture}[remember picture, overlay]
\node[anchor=south east, xshift=-1cm, yshift=1cm] at (current page.south east) {
  \begin{tikzpicture}[scale=1.5] % Increased scale for larger decorative element
    \draw[red, very thick] (0,0) circle (0.5cm);
    \draw[red, very thick] (0,0) -- (0.5,0.5);
    \draw[red, very thick] (0,0) -- (-0.5,-0.5);
    \draw[red, very thick] (0,0) -- (0.5,-0.5);
    \draw[red, very thick] (0,0) -- (-0.5,0.5);
  \end{tikzpicture}
};
\end{tikzpicture}

\title{\underline{{Assignment-1 by Shreyasi}}}
\author{SHREYASI MAHATO}
\date{July 27, 2004}


\maketitle

\begin{figure}[h!]
\centering
    \includegraphics[width=0.25\linewidth]{SHREYASI.jpg}
    \caption{SHREYASI}
    \label{fig:Shreyasi}

\end{figure}
\section{Introduction}

\subsection{About Me}
I am Shreyasi Mahato from Purulia,West Bengal. I am 20 years old.I completed my school from A.G.P.N Convent and E.R school. Currently i am studying in Maulana Abul Kalam Azad University of Technology, in Department Of Computer Application.


I’m a 20-year-old with a passion for exploring life’s many opportunities and challenges. At this exciting stage in my journey, I’m focused on discovering who I am and what I want to achieve. Whether it’s through my studies, hobbies, or spending time with friends and family, I approach life with curiosity and enthusiasm.

I’m driven by a desire to learn and grow, constantly seeking out new experiences that help me better understand the world and my place in it. I value meaningful connections with others and believe in the importance of kindness, empathy, and open-mindedness. I’m someone who enjoys deep conversations, creative pursuits, and finding joy in the little things.

As I navigate my twenties, I’m excited about the endless possibilities that the future holds. I’m determined to make the most of this time by setting goals, pursuing my passions, and staying true to myself. With a positive outlook and a strong sense of purpose, I’m ready to embrace whatever comes next.


\subsection{Interests \& Hobbies}
\begin{itemize}
    \item My hobby is doing Calligraphy.
    \item I mainly interested in Drawing and Painting
   
\end{itemize}

\subsection{Favorite Quotations}
\begin{enumerate}
    \item \textit{ "The sky is not the limit, its just the beginning."-Shanovia Lumpkin}
    \item \textit{"Flying is learning how to throw yourself at the ground and miss" -Douglas Adams"}
\end{enumerate}

\section{Mathematics}

\subsection{Mathematics and Me}
Reflecting on My Experiences with Mathematics

Mathematics has always captivated me with its power to describe the world through logic and precision. I love how it challenges me to think critically and creatively, and there’s a deep satisfaction in solving complex problems where every step leads to a definitive answer.

This year, I’ve particularly enjoyed delving into [mention specific topics, like algebra, geometry, calculus, etc.], which have broadened my perspective and deepened my appreciation for the subject. However, the journey hasn’t been without its difficulties. I found [mention the most challenging aspect, like "grappling with abstract concepts" or "mastering certain formulas"] to be especially challenging. Yet, overcoming these obstacles has been incredibly rewarding, pushing me to grow both academically and personally.

Looking ahead, I’m excited to continue my mathematical journey, with a particular interest in exploring [mention any specific interests, like applied mathematics, statistics, etc.]. I believe that mathematics is not just a subject, but a powerful tool that will be invaluable in whatever path I choose to pursue.

\subsection{Mathematical Notation}
Throughout my studies, I have frequently used mathematical notation to express complex ideas concisely. For example, I often work with quadratic equations and derivatives, using them to solve problems in physics and engineering. Mastering this notation allows me to communicate mathematical concepts clearly and effectively.


\subsubsection{Superscripts, subscripts, and Greek letters}
\begin{align*}
    &\text{(a)} & 2345 \\
    &\text{(b)} & 23_{45} \\
    &\text{(c)} & 2345^{5} \\
    &\text{(d)} & 2345_{\pi} \\
    &\text{(e)} & \cos \theta \\
    &\text{(f)} & \tan^{-1}(2.345) \\
    &\text{(g)} & \log_{23} 45 \\
    &\text{(h)} & \ln 2345 \\
    &\text{(i)} & e^{2.345} \\
    &\text{(j)} & 0 < x \leq 2345 \\
    &\text{(k)} & y \geq 2345 \\
\end{align*}

\subsubsection{Roots, fractions, and displaystyle}
\begin{align*}
    &\text{(a)} & \sqrt{2345} \\
    &\text{(b)} & \sqrt[23]{45} \\
    &\text{(c)} & \frac{23}{45} \quad \text{displaystyle:} \quad \frac{\displaystyle 23}{\displaystyle 45} \\
    &\text{(d)} & \frac{2}{3 + 4} \quad \text{displaystyle:} \quad \frac{\displaystyle 2}{\displaystyle 3 + 4} \\
    &\text{(e)} & \sqrt{\frac{23}{45}} \quad \text{displaystyle:} \quad \sqrt{\frac{\displaystyle 23}{\displaystyle 45}} \\
\end{align*}

\subsubsection{Delimiters}
\[
    \text{(a)} \quad \left( 2 + \frac{3}{45} \right)
\]
\[
    \text{(b)} \quad \left| \frac{23 - 4}{5} \right|
\]

\subsubsection{Tables and equation arrays}
\begin{center}
    \begin{tabular}{c c c c c}
        $x$ & 1 & 2 & 3 & 4 \\
        $f(x)$ & 2 & 3 & 4 & 5 \\
    \end{tabular}
\end{center}

\begin{align}
    &\text{(a)} & 2 + 3 - 4 \times 5 = x \\
    &\text{(b)} & 1 + 9 - 14 = x \\
    &\text{(c)} & 10 - 14 = x \\
    &\text{(d)} & x = -4
\end{align}

\subsubsection{Functions \& Formulas}
\begin{align*}
    &\text{(a)} & x = \frac{-b \pm \sqrt{b^2 - 4ac}}{2a} \\
    &\text{(b)} & f(x) = \left( x + \frac{2}{3} \right)^2 - \frac{4}{5} \\
    &\text{(c)} & \lim_{h \to 0} \frac{f(x+h) - f(x)}{h} = f'(x) \\
    &\text{(d)} & [f(g(x))]' = f'(g(x)) \cdot g'(x) \\
    &\text{(e)} & \frac{d^2 y}{dx^2} = f''(x) \\
    &\text{(f)} & \int \sec^2 x \, dx = \tan x + C \\
    &\text{(g)} & \int e^{2x} \, dx = \frac{1}{2} e^{2x} + C \\
    &\text{(h)} & \int_{a}^{b} f'(x) \, dx = f(b) - f(a) \\
    &\text{(i)} & \frac{d}{dx} \int_{a}^{g(x)} f(t) \, dt = f(g(x)) \cdot g'(x) \\
    &\text{(j)} & y_1 = y_0 + hF(x_0, y_0) \\
    &\text{(k)} & a_n = \{2345, \frac{2345}{2}, \frac{2345}{4}, \ldots, \frac{2345}{2^n}\} \\
    &\text{(l)} & S_n = \sum_{n=1}^{\infty} \frac{2345}{2^n} \\
    &\text{(m)} & \sum_{n=0}^{\infty} \frac{f^{(n)}(c)}{n!} (x-c)^n \\
    &\text{(n)} & \mathbf{v}(t) = x'(t) \mathbf{i} + y'(t) \mathbf{j} = \left\langle \frac{dx}{dt}, \frac{dy}{dt} \right\rangle \\
    &\text{(o)} & A = \frac{1}{2} \int_{\alpha}^{\beta} r^2 \, d\theta \\
\end{align*}
\newpage

\begin{center}
    \title\underline{{Assignment-2 by Rajib Roy}}
    \vspace{1cm}
    
    \textbf{Curriculum Vitae}

\author{Rajib Roy}
\end{center}

\date{}
\begin{wrapfigure}[14]{r}{0.19\textwidth}
    \vspace{-50pt}
    \begin{center}
    \includegraphics[width=\linewidth]{PFP.jpg}
    \end{center}
    \vspace{-10pt}
\end{wrapfigure}
\maketitle
\section*{Contact Information}
\begin{tabbing}
\hspace{3in} \= \kill
\textbf{Rajib Roy} \\                                               
Phone: +91 7029151044 \\                                                      
Email: rajib4240776@gmail.com \\
\end{tabbing}

\section*{Objective Statement}
I’m Rajib Roy, a second-semester BCA student at MAKAUT, W.B. I have a strong passion for videography, where I love capturing and creating visual stories. Along with this, I enjoy coding, finding it exciting to build and solve problems through programming. 
\section*{Education}
\textbf{Maulana Abul Kalam Azad University of Technology} \\
Bachelor of Computer Applications \\
Expected Graduation: 2027

\section*{Relevant Coursework}
\begin{itemize}[left=0pt]
    \item Programming in C 
    \item Web Development
\end{itemize}

\section*{Skills}

\textbf{Technical Skills:}
\begin{itemize}[left=0pt]
    \item \textbf{Programming Languages:} C, C++, Python, Java
    \item \textbf{Web Development:} HTML, CSS, 
    \item \textbf{Software:} MS Office, Visual Studio, Git
\end{itemize}

\textbf{Soft Skills:}
\begin{itemize}[left=0pt]
    \item Teamwork
    \item Problem-Solving
    \item Time Management
    \item Adaptability
\end{itemize}

\section*{Extracurricular Activities}
\begin{itemize}[left=0pt]
    \item \textbf{Student Clubs:} Active participant in relevant student clubs.
    \item \textbf{Volunteering:} Engaged in community service or volunteering activities.
    \item \textbf{Events:} Involved in organizing university events or workshops.
\end{itemize}

\begin{center}
    \textbf{Lab Entry By}\\
    Priyanshu
\end{center}

\subsection{Objective}
The objective of this lab is to develop a basic calculator program using the C programming language. The calculator will perform simple arithmetic operations like addition, subtraction, multiplication, and division based on user input.

\subsection{Program Overview}
The calculator program is designed to:
\begin{itemize}
    \item Accept two numbers from the user.
    \item Prompt the user to select an arithmetic operation (Addition, Subtraction, Multiplication, Division).
    \item Perform the selected operation.
    \item Display the result of the operation to the user.
\end{itemize}

The program includes error handling to manage division by zero and other invalid inputs.

\subsection{Code Implementation}
The following is the C code for the calculator program:

\begin{verbatim}
#include <stdio.h>

int main() {
    char operator;
    double num1, num2, result;

    printf("Enter an operator (+, -, *, /): ");
    scanf("%c", &operator);

    printf("Enter two operands: ");
    scanf("%lf %lf", &num1, &num2);

    switch(operator) {
        case '+':
            result = num1 + num2;
            break;
        case '-':
            result = num1 - num2;
            break;
        case '*':
            result = num1 * num2;
            break;
        case '/':
            if (num2 != 0)
                result = num1 / num2;
            else {
                printf("Error! Division by zero.\n");
                return -1;
            }
            break;
        default:
            printf("Error! Operator is not correct\n");
            return -1;
    }

    printf("Result: %.2lf\n", result);
    return 0;
}
\end{verbatim}

\subsection{Compiling and Running the Program}
To compile and run the calculator program:
\begin{enumerate}
    \item Open a terminal or command prompt.
    \item Navigate to the directory where the C file is located.
    \item Compile the program using a C compiler (e.g., GCC):
    \begin{verbatim}
    gcc calculator.c -o calculator
    \end{verbatim}
    \item Run the compiled program:
    \begin{verbatim}
    ./calculator
    \end{verbatim}
\end{enumerate}

\subsection{Adding the Calculator Program to GitHub Repository}
To add this calculator program to a GitHub repository, follow these steps:

\subsubsection{Step 1: Initialize a Local Git Repository}
\begin{enumerate}
    \item Open the terminal and navigate to the directory where your \texttt{calculator.c} file is located.
    \item If you haven't already, initialize a Git repository in that directory:
    \begin{verbatim}
    git init
    \end{verbatim}
    This command creates a new Git repository in the current directory.
\end{enumerate}

\subsubsection{Step 2: Add the File to the Repository}
\begin{enumerate}
    \item Add the \texttt{calculator.c} file to the staging area:
    \begin{verbatim}
    git add calculator.c
    \end{verbatim}
    This command stages the file, indicating that you want to include it in the next commit.
\end{enumerate}

\subsubsection{Step 3: Commit the Changes}
\begin{enumerate}
    \item Commit the file to the repository with a meaningful message:
    \begin{verbatim}
    git commit -m "Add calculator program in C"
    \end{verbatim}
\end{enumerate}

\subsubsection{Step 4: Push the Changes to GitHub}
\begin{enumerate}
    \item Link your local repository to a remote GitHub repository:
    \begin{verbatim}
    git remote add origin https://github.com/yourusername/your-repo-name.git
    \end{verbatim}
    \item Push the changes to the GitHub repository:
    \begin{verbatim}
    git push -u origin master
    \end{verbatim}
\end{enumerate}

\subsubsection{Step 5: Verify the Upload}
\begin{enumerate}
    \item Go to your GitHub repository URL in a web browser.
    \item Verify that the \texttt{calculator.c} file is listed and accessible in the repository.
\end{enumerate}
\newpage

\begin{center}
    \LARGE{\textbf{4. \underline{How to create Matrix in LaTeX}}}
\end{center}
\vspace{0.2cm}
 \begin{center}
     {\underline{Lab Assignment by Mouteer}}
\end{center}
\LARGE{To create a matrix in LaTeX, we use the \texttt{\textbf{amsmath}} package, which provides various environments for displaying matrices. Here’s a basic guide on how to create a matrix and the different options available:}
\vspace{0.1cm}
\begin{center}
\begin{itemize}
    \item{\underline{\textbf{Include the \texttt{amsmath} Package}}}
\end{itemize} 
\end{center}
First, ensure you have the \texttt{amsmath} package included in your LaTeX document preamble. Add the following line:
\begin{verbatim}
\usepackage{amsmath}
\end{verbatim}
\begin{center}
    \begin{itemize}
        \item {\textbf{\underline{Matrix Environments}}}
    \end{itemize}
\end{center}
\LARGE{
There are several environments for creating matrices, depending on how you want them formatted. Here are the most common ones:}

\begin{itemize}
  \item \texttt{matrix}: A basic matrix without brackets.
  \item \texttt{bmatrix}: A matrix with square brackets.
  \item \texttt{pmatrix}: A matrix with parentheses.
  \item \texttt{vmatrix}: A matrix with vertical bars.
  \item \texttt{Vmatrix}: A matrix with double vertical bars.
\end{itemize}
\begin{center}
        \item {\textbf{\underline{\LARGE{Syntax}}}}
\end{center}

\LARGE{The syntax for creating a matrix is similar across different environments. Use the \texttt{\textbackslash begin\{environment\}} and \texttt{\textbackslash end\{environment\}} commands to enclose the matrix content. Separate the elements in each row with \texttt{\&} and 

end each row with \texttt{\\}.

Here’s an example of a 2x2 matrix in each environment:}
\begin{verbatim}
[ 
$ \begin{matrix}
    1 & 0 \\
    0 & 1 
\end{matrix}
$ ]

\begin{bmatrix}   [
 $   \begin{matrix}
    1 & 0\\
    0 & 1
    \end{matrix}
 $   ]

a_{11} & a_{12} \\
a_{21} & a_{22}
\end{bmatrix}

\begin{pmatrix}
a_{11} & a_{12} \\
a_{21} & a_{22}
\end{pmatrix}

\begin{vmatrix}
a_{11} & a_{12} \\
a_{21} & a_{22}
\end{vmatrix}

\begin{Vmatrix}
a_{11} & a_{12} \\
a_{21} & a_{22}
\end{Vmatrix}
\end{verbatim}

\begin{center}
\section*{\textbf{\underline{\LARGE{Explanation}}}}
\end{center}
\begin{itemize}
  \item \texttt{\textbackslash begin\{matrix\} ... \textbackslash end\{matrix\}}: Creates a matrix with no surrounding brackets.
  \item \texttt{\textbackslash begin\{bmatrix\} ... \textbackslash end\{bmatrix\}}: Surrounds the matrix with square brackets.
  \item \texttt{\textbackslash begin\{pmatrix\} ... \textbackslash end\{pmatrix\}}: Surrounds the matrix with parentheses.
  \item \texttt{\textbackslash begin\{vmatrix\} ... \textbackslash end\{vmatrix\}}: Surrounds the matrix with single vertical bars, often used to denote determinants.
  \item \texttt{\textbackslash begin\{Vmatrix\} ... \textbackslash end\{Vmatrix\}}: Surrounds the matrix with double vertical bars, also used for determinants or norms.
\end{itemize}

\newpage
\textbf{Lab Assignment 5: Mind Reader App}

\begin{center}
    \textbf{Entry by Abhirup}
\end{center}


\textbf{Task:} Your professor created a mind reader application and wants you to try it out. After running the program, you found the submit button looks dull. You renamed it "Chin Tapak Dum Dum," but the button became disproportionate. Your task is to fix the button issue and create a pull request with the solution.

\subsection*{Procedure}
\begin{enumerate}[label=\arabic*.]
    \item \textbf{Clone the Repository:}
    \begin{itemize}
        \item Open GitHub Desktop or use the terminal to clone the repository: \url{https://github.com/GeekAyan/STT}
        \item Run the command: \texttt{git clone https://github.com/GeekAyan/STT.git}
        \item Navigate to the project directory.
    \end{itemize}
    \item \textbf{Run the Application:}
    \begin{itemize}
        \item Follow the instructions provided in the \texttt{README.md} file to set up and run the mind reader application using your preferred Integrated Development Environment (IDE).
        \item Observe the application's user interface, particularly the submit button.
    \end{itemize}
    \item \textbf{Identify and Rename the Button:}
    \begin{itemize}
        \item Locate the submit button code in the application’s source files.
        \item Rename the button text to "Chin Tapak Dum Dum."
        \item Notice that the button has become disproportionate due to the increased text length.
    \end{itemize}
    \item \textbf{Fix the Button Size:}
    \begin{itemize}
        \item Analyze the layout code that controls the button's appearance.
        \item Adjust the width and height properties, or use appropriate CSS/JavaFX adjustments to make the button proportionate.
        \item Test the application to ensure the button now displays correctly and does not affect other UI elements.
    \end{itemize}
    \item \textbf{Commit and Push the Changes:}
    \begin{itemize}
        \item Stage the modified files with: \texttt{git add .}
        \item Commit the changes with a descriptive message: \texttt{git commit -m "Fix button size after renaming to 'Chin Tapak Dum Dum'"}
        \item Push the changes to your forked repository on GitHub.
    \end{itemize}
    \item \textbf{Create a Pull Request:}
    \begin{itemize}
        \item Go to your GitHub repository and click on "Compare and pull request."
        \item Write a brief description of the changes made and submit the pull request to the original repository.
    \end{itemize}
    \item \textbf{Review and Merge:}
    \begin{itemize}
        \item Wait for the repository owner to review your pull request.
        \item If accepted, your changes will be merged into the main project.
    \end{itemize}
\end{enumerate}

\end{document}


\end{document}
